% Claim proof file for row claim_3_14 -- "AddingInterventionNodes".
%
% Proof lifted from lecture-notes/lecture_notes/graphs.tex:845--875 (the
% \Claude{...} block immediately following the claim). Light edits:
% \refrow{def_3_10} / \refrow{def_3_13} cross-links on first use,
% \medskip\noindent\emph{...} per-component headers mirroring the sibling
% claim_3_11_proof_DisjointHardInterventions.tex, a \checkmark at the end
% of each component-agree paragraph, the LN's terse "Output nodes ($V$)
% and bidirected edges ($L$) are unchanged" telescoped phrase split into
% its own $V$-paragraph and $L$-paragraph in chain 1, an explicit
% constructor-disjointness argument behind the LN's load-bearing
% sub-identities $W_2 \sm J_1 = W_2 \sm J$ (chain 1) and
% $W_1 \sm (J \cup W_2) = W_1 \sm J$ (chain 2), and a closing one-liner.
% The \Claude{...} wrapper from the LN source is dropped.
%
% Compiles standalone via the `subfiles` package.

\documentclass[main]{subfiles}

\begin{document}
% ref: claim_3_14 (proof, with statement restated for readability)
% title: AddingInterventionNodes
\def\rowref{claim\_3\_14}
\phantomsection\label{claim_3_14}
%
% Cross-references: see claim_<ref>_statement_<title>.tex in this folder
% for the corresponding statement.

% --- statement (restated) ---------------------------------------------------
\begin{Lem}[Adding intervention nodes commutes with disjoint hard interventions]
    Let $G=(J,V,E,L)$ be a CDMG and $W_1, W_2 \ins J \cup V$ two disjoint subsets of nodes from $G$.
      Then we have:
      \[ \lp G_{\doit(I_{W_1})} \rp_{\doit(I_{W_2})} = \lp G_{\doit(I_{W_2})} \rp_{\doit(I_{W_1})} 
      =  G_{\doit(I_{W_1 \cup W_2})}. \]
      We also have:
      \[ \lp G_{\doit(I_{W_1})} \rp_{\doit(W_2)} = \lp G_{\doit(W_2)} \rp_{\doit(I_{W_1})} =  G_{\doit(I_{W_1}, W_2)}. \]
   % So there is not ambiguity in the following notations and identifications:
   % \[ G(V|\doit(J\cup I_{W_1} \cup I_{W_2})) = G(V|\doit(J\cup I_{W_1 \cup W_2})),
   % \qquad G(V\sm W_2|\doit(W_2 \cup J\cup I_{W_1})).\]
\end{Lem}

% --- proof ------------------------------------------------------------------
\begin{proof}
\medskip\noindent\emph{Preamble.}
Let $G = (J, V, E, L)$ be a CDMG and $W_1, W_2 \ins J \cup V$ two disjoint
subsets of nodes of $G$. The lemma bundles two chained equalities, one for
the iterated extension by intervention nodes \refrow{def_3_13} (first chain)
and one for the mixed extend / hard-intervene composite~\refrow{def_3_10}
(second chain). In each case we present every CDMG as a 4-tuple
$(J_\ast, V_\ast, E_\ast, L_\ast)$ and verify that the two iterates agree
componentwise, in turn matching the LN's third term where one is named.

The load-bearing observation, used in both chains, is that the fresh
intervention nodes minted by~\refrow{def_3_13} are \emph{constructor-disjoint}
from the original node universe $J \cup V$: for every $w \in W \sm J$ the
new input vertex $I_w$ is, by construction, a fresh symbol that equals
neither any $v \in J \cup V$ nor any $I_{w'}$ with $w \neq w'$. Two
immediate consequences will be used repeatedly:
\begin{itemize}
  \item \emph{(no collision with fresh intervention nodes, chain 1).}
        Writing $J_1 := J \dcup \lC I_w \mid w \in W_1 \sm J \rC$ for the
        input set of $G_{\doit(I_{W_1})}$, we have
        $W_2 \sm J_1 = W_2 \sm J$. Indeed any $w \in W_2 \ins J \cup V$ is
        original (so $w \neq I_{w'}$ for every $w'$), and $w \notin W_1$ by
        $W_1 \cap W_2 = \emptyset$, hence $w \notin W_1 \sm J$, hence
        $w \notin J_1$ iff $w \notin J$.
  \item \emph{(no fresh edges into $W_2$, chain 2).} Since each fresh edge
        $I_w \tuh w$ minted by~\refrow{def_3_13} has head $w \in W_1$ (the
        target of the extension), and $W_1 \cap W_2 = \emptyset$, no fresh
        edge has its head in $W_2$. Equivalently, $W_1 \sm (J \cup W_2) =
        W_1 \sm J$: a $w \in W_1$ is automatically outside $W_2$ by
        disjointness, so the extra $\notin W_2$ cut on $W_1$ is vacuous.
\end{itemize}

\medskip\noindent\emph{First chain:
$\lp G_{\doit(I_{W_1})} \rp_{\doit(I_{W_2})}
   = \lp G_{\doit(I_{W_2})} \rp_{\doit(I_{W_1})}
   = G_{\doit(I_{W_1 \cup W_2})}$.}

We first prove $\lp G_{\doit(I_{W_1})} \rp_{\doit(I_{W_2})}
= G_{\doit(I_{W_1 \cup W_2})}$ by checking the four components $(J, V, E, L)$.

\medskip\noindent\emph{(a) Input nodes ($J$).}
By~\refrow{def_3_13}, the inner extension gives
\[
  J_{G_{\doit(I_{W_1})}}
  \;=\; J \dcup \lC I_w \mid w \in W_1 \sm J \rC
  \;=:\; J_1.
\]
Applying~\refrow{def_3_13} again to extend by $W_2$ on top of $G_{\doit(I_{W_1})}$,
\[
  J_{\lp G_{\doit(I_{W_1})}\rp_{\doit(I_{W_2})}}
  \;=\; J_1 \dcup \lC I_w \mid w \in W_2 \sm J_1 \rC.
\]
By the preamble's no-collision identity, $W_2 \sm J_1 = W_2 \sm J$, so
\[
  J_{\lp G_{\doit(I_{W_1})}\rp_{\doit(I_{W_2})}}
  \;=\; J \dcup \lC I_w \mid w \in W_1 \sm J \rC \dcup \lC I_w \mid w \in W_2 \sm J \rC
  \;=\; J \dcup \lC I_w \mid w \in (W_1 \cup W_2) \sm J \rC,
\]
where the last equality unions the two indexing sets (the disjointness
$(W_1 \sm J) \cap (W_2 \sm J) = \emptyset$ follows from $W_1 \cap W_2
= \emptyset$). This is exactly $J_{G_{\doit(I_{W_1 \cup W_2})}}$ produced
by a single application of~\refrow{def_3_13} on $W_1 \cup W_2$. \checkmark

\medskip\noindent\emph{(b) Output nodes ($V$).}
By~\refrow{def_3_13}, every application of $\doit(I_\cdot)$ preserves the
output set verbatim. Iterating twice on the left and once on the right,
\[
  V_{\lp G_{\doit(I_{W_1})}\rp_{\doit(I_{W_2})}}
  \;=\; V_{G_{\doit(I_{W_1})}}
  \;=\; V
  \;=\; V_{G_{\doit(I_{W_1 \cup W_2})}}. \quad\checkmark
\]

\medskip\noindent\emph{(c) Directed edges ($E$).}
By~\refrow{def_3_13},
\[
  E_{G_{\doit(I_{W_1})}}
  \;=\; E \dcup \lC I_w \tuh w \mid w \in W_1 \sm J \rC
  \;=:\; E_1,
\]
and a second application on $W_2$ gives
\[
  E_{\lp G_{\doit(I_{W_1})}\rp_{\doit(I_{W_2})}}
  \;=\; E_1 \dcup \lC I_w \tuh w \mid w \in W_2 \sm J_1 \rC.
\]
Using $W_2 \sm J_1 = W_2 \sm J$ from the preamble and unioning the indexing
sets (disjoint by $W_1 \cap W_2 = \emptyset$),
\[
  E_{\lp G_{\doit(I_{W_1})}\rp_{\doit(I_{W_2})}}
  \;=\; E \dcup \lC I_w \tuh w \mid w \in (W_1 \cup W_2) \sm J \rC
  \;=\; E_{G_{\doit(I_{W_1 \cup W_2})}}. \quad\checkmark
\]

\medskip\noindent\emph{(d) Bidirected edges ($L$).}
By~\refrow{def_3_13}, every application of $\doit(I_\cdot)$ preserves the
bidirected-edge set verbatim. Iterating twice on the left and once on the right,
\[
  L_{\lp G_{\doit(I_{W_1})}\rp_{\doit(I_{W_2})}}
  \;=\; L_{G_{\doit(I_{W_1})}}
  \;=\; L
  \;=\; L_{G_{\doit(I_{W_1 \cup W_2})}}. \quad\checkmark
\]

All four components of $\lp G_{\doit(I_{W_1})} \rp_{\doit(I_{W_2})}$ agree
with those of $G_{\doit(I_{W_1 \cup W_2})}$, so the two CDMGs are equal.

\medskip\noindent\emph{By symmetry, $\lp G_{\doit(I_{W_2})} \rp_{\doit(I_{W_1})}
= G_{\doit(I_{W_1 \cup W_2})}$.}
The argument in (a)--(d) is symmetric in $W_1$ and $W_2$: the hypothesis
$W_1 \cap W_2 = \emptyset$ is preserved under swapping (it equals
$W_2 \cap W_1 = \emptyset$), the no-collision identity becomes
$W_1 \sm J_2 = W_1 \sm J$ where $J_2 := J \dcup \lC I_w \mid w \in W_2 \sm J \rC$
(same proof, with $1$ and $2$ exchanged), and the unioned indexing set
$W_2 \cup W_1$ in (a) and (c) equals $W_1 \cup W_2$ by commutativity of
$\cup$. Components (b) and (d) hold verbatim. Hence the two iterates of
chain 1 are both equal to $G_{\doit(I_{W_1 \cup W_2})}$, and the three-term
chain closes.

\medskip\noindent\emph{Second chain:
$\lp G_{\doit(I_{W_1})} \rp_{\doit(W_2)}
   = \lp G_{\doit(W_2)} \rp_{\doit(I_{W_1})}
   = G_{\doit(I_{W_1}, W_2)}$.}

A remark on the third term: as the LN's commented-out lines
(\texttt{graphs.tex}~840--842) make explicit, $G_{\doit(I_{W_1}, W_2)}$
is \emph{notation for the unambiguous value of either iterate}, not a
separate construction. The content of the chain is therefore the single
equation between the two iterates; verifying it componentwise establishes
the entire three-term chain at once.

Write $J_1 := J \dcup \lC I_w \mid w \in W_1 \sm J \rC$ and
$E_1 := E \dcup \lC I_w \tuh w \mid w \in W_1 \sm J \rC$ for the input
set and edge set of $G_{\doit(I_{W_1})}$ (so that $G_{\doit(I_{W_1})} =
(J_1, V, E_1, L)$ by~\refrow{def_3_13}).

\medskip\noindent\emph{(a) Input nodes ($J$).}
Starting from $G_{\doit(I_{W_1})}$ and applying $\doit(W_2)$
via~\refrow{def_3_10}, the hard intervention sends the input set to
$J_1 \cup W_2$:
\[
  J_{\lp G_{\doit(I_{W_1})}\rp_{\doit(W_2)}}
  \;=\; J_1 \cup W_2
  \;=\; \lp J \dcup \lC I_w \mid w \in W_1 \sm J \rC \rp \cup W_2.
\]
For the other order, by~\refrow{def_3_10} the inner hard intervention gives
$J_{G_{\doit(W_2)}} = J \cup W_2$, and~\refrow{def_3_13} applied on top with
$W_1$ extends by $\lC I_w \mid w \in W_1 \sm (J \cup W_2) \rC$:
\[
  J_{\lp G_{\doit(W_2)}\rp_{\doit(I_{W_1})}}
  \;=\; (J \cup W_2) \dcup \lC I_w \mid w \in W_1 \sm (J \cup W_2) \rC.
\]
The preamble's no-fresh-edge-into-$W_2$ identity rewrites the indexing
set: $W_1 \sm (J \cup W_2) = W_1 \sm J$, since every $w \in W_1$ is
already outside $W_2$ by $W_1 \cap W_2 = \emptyset$. So
\[
  J_{\lp G_{\doit(W_2)}\rp_{\doit(I_{W_1})}}
  \;=\; (J \cup W_2) \dcup \lC I_w \mid w \in W_1 \sm J \rC,
\]
which agrees with $J_{\lp G_{\doit(I_{W_1})}\rp_{\doit(W_2)}}$ above up
to commutativity of $\cup$ on the original-input piece $J \cup W_2$. \checkmark

\medskip\noindent\emph{(b) Output nodes ($V$).}
Left: by~\refrow{def_3_13} the inner $\doit(I_{W_1})$ preserves $V$, and
by~\refrow{def_3_10} the outer $\doit(W_2)$ deletes $W_2$ from the output
set, giving $V_{\lp G_{\doit(I_{W_1})}\rp_{\doit(W_2)}} = V \sm W_2$.
Right: by~\refrow{def_3_10} the inner $\doit(W_2)$ deletes $W_2$ from
the output set, $V_{G_{\doit(W_2)}} = V \sm W_2$; and by~\refrow{def_3_13}
the outer $\doit(I_{W_1})$ preserves the output set, giving
$V_{\lp G_{\doit(W_2)}\rp_{\doit(I_{W_1})}} = V \sm W_2$. \checkmark

\medskip\noindent\emph{(c) Directed edges ($E$).}
This is the load-bearing case (LN: \texttt{graphs.tex}~866--871). Left:
starting from $E_{G_{\doit(I_{W_1})}} = E_1$ and applying $\doit(W_2)$
via~\refrow{def_3_10}.iii, the hard intervention deletes every edge with
head in $W_2$:
\[
  E_{\lp G_{\doit(I_{W_1})}\rp_{\doit(W_2)}}
  \;=\; E_1 \sm \lC v \tuh w \mid w \in W_2 \rC.
\]
Of the edges in $E_1 = E \dcup \lC I_w \tuh w \mid w \in W_1 \sm J \rC$,
the fresh edges $\lC I_w \tuh w \mid w \in W_1 \sm J \rC$ all have head
in $W_1$, which is disjoint from $W_2$ by hypothesis, hence none of them
is deleted; only original $E$-edges with head in $W_2$ are removed. So
\[
  E_{\lp G_{\doit(I_{W_1})}\rp_{\doit(W_2)}}
  \;=\; \lp E \sm \lC v \tuh w \mid w \in W_2 \rC \rp \dcup \lC I_w \tuh w \mid w \in W_1 \sm J \rC.
\]
Right: by~\refrow{def_3_10}.iii, $E_{G_{\doit(W_2)}} = E \sm \lC v \tuh w
\mid w \in W_2 \rC$. Applying~\refrow{def_3_13} to extend by $W_1$ then
appends the fresh edges with indexing set $W_1 \sm (J \cup W_2) = W_1 \sm J$
(from the preamble):
\[
  E_{\lp G_{\doit(W_2)}\rp_{\doit(I_{W_1})}}
  \;=\; \lp E \sm \lC v \tuh w \mid w \in W_2 \rC \rp \dcup \lC I_w \tuh w \mid w \in W_1 \sm J \rC.
\]
Same expression as the left. \checkmark

\medskip\noindent\emph{(d) Bidirected edges ($L$).}
Left: by~\refrow{def_3_13} the inner $\doit(I_{W_1})$ preserves the
bidirected-edge set, giving $L_{G_{\doit(I_{W_1})}} = L$; and
by~\refrow{def_3_10}.iv the outer $\doit(W_2)$ removes every bidirected
edge incident to $W_2$, giving
$L_{\lp G_{\doit(I_{W_1})}\rp_{\doit(W_2)}} = L \sm \lC v \huh w \mid w \in W_2 \rC$.
Right: by~\refrow{def_3_10}.iv, $L_{G_{\doit(W_2)}} = L \sm \lC v \huh w
\mid w \in W_2 \rC$, and the outer $\doit(I_{W_1})$ preserves the
bidirected-edge set by~\refrow{def_3_13}, giving
$L_{\lp G_{\doit(W_2)}\rp_{\doit(I_{W_1})}} = L \sm \lC v \huh w \mid w \in W_2 \rC$.
\checkmark

All four components agree, completing the second chain. Both chains have
been verified componentwise; this completes the proof.
\end{proof}

\end{document}
